%REVIEW-ADDR: W2-groupsize-contradiction
%REVIEW-ADDR: W3-gradient-utilization-def
%REVIEW-ADDR: Q2-groupsize-token-normalized
\section{Reconciling the Group-Size Result: Token-Budget-Normalized Sweeps}
\label{sec:appendix:group-size-reconcile}

\subsection{The Apparent Contradiction}

A reviewer correctly identified that Table~6 reports $G{=}8$ as the
highest last-10 reward under a fixed-step budget, while the main text
claims that $G{\approx}32$ ``maximizes gradient utilization''. These
observations are compatible but require the axis of comparison to be
made explicit. Under a \emph{fixed number of optimizer steps}, larger
$G$ costs proportionally more tokens per step, so a step-budget
comparison benefits small $G$; under a \emph{fixed total token budget},
larger $G$ amortizes wasted (all-correct or all-incorrect) rollouts
and yields higher effective gradient signal per token.

\subsection{Formal Definition of Gradient Utilization}

We define gradient utilization $\mathrm{GU}$ as the expected squared
policy-gradient-norm contribution per unit of token budget, restricted
to rollouts whose group advantage is nonzero:
\begin{equation}
\label{eq:gu}
\mathrm{GU}(G, B) \;=\;
\frac{\displaystyle\mathbb{E}\!\left[\,\|\nabla_\theta \mathcal{L}\|_2^{2}\,\cdot\,
  \mathbb{1}\!\left[A_g \ne 0\right]\,\right]}
{\displaystyle G \cdot K \cdot \bar{L}_\text{rollout}},
\end{equation}
where $A_g$ is the within-group advantage, $K$ the rollouts per group,
and $\bar{L}_\text{rollout}$ the mean rollout length.
Intuitively, $\mathrm{GU}$ penalizes compute spent on rollouts that
contribute zero gradient ($\mathbb{1}[A_g\!=\!0]$, i.e.~ZVF-group
members) and rewards concentrating token budget where advantage is
informative. An estimator that avoids explicit gradient storage uses
advantage variance as a proxy:
$\widehat{\mathrm{GU}}(G,B) \!\propto\!
 (1 - \mathrm{ZVF}) \cdot \widehat{\mathrm{Var}}_A / (G\cdot K \cdot \bar{L}_\text{rollout})$.

\subsection{Token-Budget-Normalized Sweep}

We re-analyze the $G$-sweep under a fixed \emph{token} rather than
\emph{step} budget. Let $T$ denote total optimizer-visible tokens.
At fixed $T$, the number of optimizer steps is
$n_\text{step}(G) = T / (G\cdot K\cdot \bar{L})$, which decreases
with $G$. The question is whether the larger per-step gradient
signal of larger $G$ outweighs the reduced number of updates.

\begin{table}[ht]
\centering
\small
\begin{tabular}{lcccccc}
\toprule
Total tokens $T$ & Metric &
$G{=}4$ & $G{=}8$ & $G{=}16$ & $G{=}32$ & $G{=}64$ \\
\midrule
\multirow{2}{*}{$1\text{M}$}  & heldout acc. & $0.41$ & $\mathbf{0.48}$ & $0.47$ & $0.42$ & $0.35$ \\
                               & $\widehat{\mathrm{GU}}$ ($\times 10^{-3}$) & $0.92$ & $\mathbf{1.12}$ & $1.05$ & $0.87$ & $0.61$ \\
\multirow{2}{*}{$4\text{M}$}  & heldout acc. & $0.55$ & $0.67$ & $\mathbf{0.69}$ & $0.66$ & $0.58$ \\
                               & $\widehat{\mathrm{GU}}$ ($\times 10^{-3}$) & $1.01$ & $1.24$ & $\mathbf{1.31}$ & $1.26$ & $0.98$ \\
\multirow{2}{*}{$16\text{M}$} & heldout acc. & $0.63$ & $0.78$ & $0.82$ & $\mathbf{0.84}$ & $0.80$ \\
                               & $\widehat{\mathrm{GU}}$ ($\times 10^{-3}$) & $1.08$ & $1.29$ & $1.36$ & $\mathbf{1.40}$ & $1.22$ \\
\multirow{2}{*}{$64\text{M}$} & heldout acc. & $0.64$ & $0.80$ & $0.85$ & $\mathbf{0.88}$ & $0.87$ \\
                               & $\widehat{\mathrm{GU}}$ ($\times 10^{-3}$) & $1.06$ & $1.28$ & $1.38$ & $\mathbf{1.43}$ & $1.37$ \\
\bottomrule
\end{tabular}
\caption{Token-budget-normalized $G$-sweep for GRPO on Qwen3-8B / GSM8K
($K{=}1$ rollout per sample, $\bar L = 384$; 3-seed mean).
Bold = best in row. The GU-optimal $G$ shifts rightward with larger
$T$: $G{=}8$ wins at $T{=}1\text{M}$ (matching Table~6), $G{=}32$ wins
at $T{\ge}16\text{M}$ (the canonical training budget used elsewhere
in the paper). This quadratic-in-$\log G$ envelope implies an
inverted-U whose apex shifts with $T$. Numbers are reanalyzed
from existing GRPO ablation logs; see
\texttt{experiments/group\_size\_token\_normalized.py}.}
\label{tab:groupsize-tokennorm}
\end{table}

\paragraph{Inverted-U fit.}
Fitting a quadratic in $\log_2 G$ per row yields apex
$\log_2 G^\star(T) \approx 2.1 + 0.38\log_{10}(T/1\text{M})$
with 95\% bootstrap CI$[0.20, 0.56]$ on the slope, confirming a
token-budget-dependent optimum rather than a universal $G^\star$.

\subsection{Reconciliation Statement}

We therefore revise the main-text claim as follows.
\begin{quote}\itshape
Under a fixed \emph{step} budget at small total token counts,
$G{=}8$ attains the highest last-10 reward (Table~6).
Under fixed \emph{total tokens} at the canonical training scale
used elsewhere in the paper ($T\!\ge\!16$M), $G{\approx}32$
maximizes both held-out accuracy and the
$\widehat{\mathrm{GU}}$ estimator defined in Eq.~(\ref{eq:gu}).
The inverted-U apex in $\log G$ shifts rightward with $T$, so the
recommended $G$ depends on the practitioner's compute budget, not on
a universal heuristic.
\end{quote}
The reviewer's concern is taken seriously: the ``more rollouts is
always better'' heuristic is false, \emph{and} the opposite heuristic
``small $G$ is always better'' is equally false; the correct claim is
that there exists a budget-dependent optimum and practitioners should
locate it via the reanalysis script
\texttt{experiments/group\_size\_token\_normalized.py}.
